\section{Round-five referee closure: canonical strategy laws, relaxed DPP, testing, and projective LAN}
\label{sec:r5-c1}

Adaptive controls are constructed on one canonical history/control space by
Ionescu--Tulcea recursion.  The sufficient finite-scale state is the complete
resolved history together with the current posterior and hierarchy law; a
slow coordinate alone is never declared Markov.

\subsection{Canonical strategy space}

Let
\[
 \Omega^{\rm can}
 =
 {\mathsf H}\times
 D([0,T],U)\times D([0,T],V)\times
 {\mathcal P}(Z)\times{\mathfrak X}
\]
carry the coordinate history \(H\), controls \(u,v\), posterior \(\rho\),
and complete hierarchy state \(G\).  On a partition
\(0=t_0<\cdots<t_m=T\), an Elliott--Kalton strategy is a Borel map from the
opponent's past controls and the resolved history to the next control.

For a block source \(Y_k^{u,v}\), define the conditional normalizer under the
law already constructed through time \(t_k\):
\[
 {\mathcal N}_k^{u,v}
 =
 \log\mathbb E^{u,v}\!\left[
 e^{\mu_\varepsilon Y_k^{u,v}}
 \mid{\mathcal F}^{\rm can}_{t_k}\right].
\]

\begin{theorem}[Pathwise-consistent controlled law]
\label{thm:r5-c1-law}
For every pair of nonanticipative block strategies there is a unique
probability law on \(\Omega^{\rm can}\) such that
\[
 \frac{d\mathbb P^{u,v}_{k+1}}
      {d\mathbb P^{u,v}_{k}}
 =
 \exp\{\mu_\varepsilon Y_k^{u,v}-{\mathcal N}_k^{u,v}\}.
\]
The density process is a positive mean-one martingale.  If two strategy pairs
agree on the realized history through \(t_k\), their laws agree on
\({\mathcal F}^{\rm can}_{t_k}\) and have the same next conditional
normalizer.  Thus opponent-dependent strategies are pathwise consistent.
\end{theorem}

\begin{proof}
Assume the law is constructed through \(t_k\).  The strategies are Borel
functions of the canonical past, so the next controls and source are
measurable.  The B2 finite-source block kernel gives a regular conditional
law and a finite conditional exponential moment.  Normalization makes the
displayed density have conditional mean one.  Ionescu--Tulcea recursion
therefore gives a unique law on the next block.  Induction gives the full law
and martingale property.  Equality of past strategy evaluations yields
equality of the conditional kernels and hence of the normalizers.
\end{proof}

\subsection{Relaxed controls and dynamic programming}

Let \({\mathcal R}_U\) and \({\mathcal R}_V\) be Young measures on time and
control, with the stable topology.  The controlled collision multiplier is
the barycentric kernel
\[
 q^{\alpha,\beta}
 =\int q^{u,v}\,\alpha_t(du)\beta_t(dv),
\]
bounded between fixed positive constants.

\begin{lemma}[Compactness and lower semicontinuity]
\label{lem:r5-c1-compact}
The relaxed strategy laws with bounded action are tight on
\(\Omega^{\rm can}\).  Every limit is nonanticipative, the controlled balance
is closed, and
\[
 \int\ell\left(\frac{d\Gamma}
 {q^{\alpha,\beta}dA_f}\right)
 q^{\alpha,\beta}dA_f
\]
is lower semicontinuous.
\end{lemma}

\begin{proof}
Compactness of the control sets gives compactness of Young measures.
B2 exponential compact containment gives tightness of density, collision
current, and hierarchy coordinates; A4 gives history tightness.  Stable
convergence preserves integrals of bounded predictable kernels and therefore
nonanticipativity.  Closedness of balance is B3's weighted result.
The entropy integrand is the perspective of the convex function \(\ell\),
hence is jointly lower semicontinuous under weak convergence and uniform
positive bounds on \(q\).
\end{proof}

\begin{theorem}[Relaxed Elliott--Kalton DPP]
\label{thm:r5-c1-dpp}
The upper and lower relaxed values satisfy the dynamic programming principle
on the state
\[
 (H_t,\rho_t,G_t).
\]
If the compact-control Isaacs equality holds, they coincide and are the
unique viscosity solution of
\[
 \partial_tU+
 \sup_\alpha\inf_\beta
 \left\{
 \langle f,v\cdot\nabla p\rangle+
 \int q^{\alpha,\beta}(e^{\Delta p}-1)dA_f+
 c(f,\alpha,\beta)
 \right\}=0
\]
on the B4 exponential-moment domain.
\end{theorem}

\begin{proof}
Theorem~\ref{thm:r5-c1-law} permits concatenation of strategy laws because
conditional normalizers are recomputed from the current canonical law.
Measurable selection holds first on finite partitions; compactness and
Lemma~\ref{lem:r5-c1-compact} pass it to relaxed limits.  The short-block B2
Laplace principle identifies the Hamiltonian.  B4 comparison extends after
taking \(\sup\inf\) because the controlled multipliers are uniformly bounded
and equicontinuous.  Isaacs equality identifies the two half-relaxed limits.
\end{proof}

\subsection{Information state and complete observation}

Let
\[
 {\mathcal F}^{\rm res}_t
 =\sigma(H_s,\rho_s,G_s: s\le t).
\]

\begin{proposition}[Sufficiency and deterministic continuation]
\label{prop:r5-c1-sufficiency}
The conditional future law under every prepared phase and admissible strategy
is a Borel function of \((H_t,\rho_t,G_t)\).  On the full-measure regular set
where the complete microscopic phase point is observed, this conditional law
is a Dirac mass at the unique Hamiltonian continuation.
\end{proposition}

\begin{proof}
The hierarchy coordinate determines the conditional microscopic law, the
history determines all predictable sources and controls, and the posterior
determines the phase mixture.  Disintegration therefore factors through the
declared state.  For a regular hard-sphere or billiard phase point,
Hamiltonian/specular dynamics has a unique continuation until the next
regular contact and then recursively thereafter; grazing and simultaneous
contacts are null.  Hence complete observation produces the stated Dirac law.
\end{proof}

\subsection{Phase testing exponent}

Let \((\mathbb P_z)_{z\in Z}\) be a compact family of prepared phases on the
same canonical history space.  For \(z_0,z\), define the path relative-entropy
rate
\[
 {\mathcal K}_T(z_0\Vert z)
 =
 \inf_{\nu:\nu\text{ has the }z_0\text{ mean statistics}}
 I_{T,z}(\nu).
\]

\begin{theorem}[Expected posterior testing exponent]
\label{thm:r5-c1-testing}
Assume a finite family of bounded history/density/contact statistics separates
\(z_0\) from the closed set \(F\subset Z\).  Then
\[
 \limsup_{\varepsilon\to0}
 \frac1{\mu_\varepsilon}
 \log\mathbb E_{z_0}\rho_T(F)
 \le
 -\inf_{z\in F}{\mathcal K}_T(z_0\Vert z)<0.
\]
The infimum uses the full B2 rate, including microcanonical conditioning and
actual-contact recovery.
\end{theorem}

\begin{proof}
Cover \(F\) by finitely many neighbourhoods on which one separating statistic
has a uniform mean gap.  Bayes' formula writes posterior mass as an integral
of likelihood ratios.  Apply the B2 full LDP upper bound under the alternative
phase and the B1 conditioned normalization.  The recovery theorem makes the
dual relative-entropy expression exact, not merely an exposed-point bound.
Compactness of \(F\) and the separating gap give a strictly positive uniform
infimum.  Integrating against the prior proves the estimate.
\end{proof}

\subsection{Projective local asymptotic normality}

Let \({\mathcal H}_0\) be the countable Hilbert source space used in B3 and
let \(h\) range over a fixed finite-dimensional subspace
\(E\subset{\mathcal H}_0\).  The finite-volume tilt is centered at the exact
B1 saddle.

\begin{theorem}[Canonical-chart LAN and posterior Gaussian tangent]
\label{thm:r5-c1-lan}
For every finite-dimensional \(E\) and bounded \(h\in E\),
\[
 \log\frac{d\mathbb P_{\Theta+h/\sqrt{\mu_\varepsilon},\varepsilon}}
               {d\mathbb P_{\Theta,\varepsilon}}
 =
 \langle h,Z_\varepsilon\rangle
 -\frac12\Sigma_\Theta(h,h)+o_{\mathbb P_\Theta}(1),
\]
uniformly on compact \(h\)-sets, where \(Z_\varepsilon\) is centered by the
finite-volume mean and converges to the B3 balanced Gaussian field.  The
likelihood ratios are uniformly integrable.  Compatible finite-dimensional
statements define projective LAN on the countable source chart.  For a prior
with positive continuous density in local coordinates, the rescaled posterior
converges to the corresponding Gaussian shift experiment.
\end{theorem}

\begin{proof}
B2 normal convergence on a larger complex finite-source ball gives a uniform
third-order Cauchy remainder for
\(h/\sqrt{\mu_\varepsilon}\).  B1's exact finite saddle removes the
finite-volume linear mismatch.  B3 gives convergence of \(Z_\varepsilon\)
and the Hessian covariance.  A slightly larger real source ball gives an
\(L^{1+\delta}\) likelihood bound, proving uniform integrability.
Compatibility under restriction of \(E\) gives the projective statement.
Taylor expansion of the prior and the LAN likelihood, followed by dominated
convergence, gives the posterior Gaussian limit.
\end{proof}

\begin{theorem}[Round-five C1 closure]
\label{thm:r5-c1-main}
Adaptive laws are defined consistently for opponent-dependent strategies,
relaxed controls satisfy a complete-history DPP, posterior testing uses the
full conditioned path rate, and LAN holds on the compatible canonical
exponential chart with exact finite-volume centering.
\end{theorem}

\begin{proof}
Combine Theorems~\ref{thm:r5-c1-law},
\ref{thm:r5-c1-dpp}, \ref{thm:r5-c1-testing}, and
\ref{thm:r5-c1-lan}.
\end{proof}
